% FORMALIZACION_NUEVA of the signed incidence of the preexisting Paley chart.
% The construction starts at exc:aw; it does not identify its flags by decree
% with the quotient of all TPK histories.
\paragraph{Construction of the carrier from six-position incidence.}
\label{ii:ico:construccion-firmada}
The matrix \(A_W\) of \eqref{exc:aw}, which previously enters
angular transport in section~\ref{sec:ii-enlace-angular-witt}, provides
an explicit antecedent of the icosahedral realization. Let
\(C=\operatorname{bal}(A_W)\), where
\(\operatorname{bal}(0)=0\), \(\operatorname{bal}(1)=1\), and
\(\operatorname{bal}(2)=-1\). Its coordinate order is
\(\mathcal U=(\infty,0,1,2,3,4)\). The integer chart satisfies
\(C=C^{\mathsf T}\), \(C^2=5I_6\), and \(\operatorname{Tr}C=0\).
Define its signed covering by
\begin{equation}
 \mathcal D_C=\mathcal U\times\{+1,-1\},\qquad
 (u,\sigma)\sim_C(v,\tau)
 \ \Longleftrightarrow\ u\ne v\ \text{ and }\ \sigma\tau=C_{uv}.
 \label{ii:ico:incidencia-firmada}
\end{equation}
Orientation doubles each coordinate; the matrix decides the incidences
between the two corresponding fibers. The involution of this covering
is \(\jmath_C(u,\sigma)=(u,-\sigma)\).

\Needspace{12\baselineskip}
\begin{samepage}
\begin{proposition}[Isometric realization of signed incidence]
\label{ii:ico:realizacion-firmada}
The operator and vectors
\begin{equation}
 E_C=\frac12\left(I_6+\frac C{\sqrt5}\right),\qquad
 \mathcal H_C=\operatorname{im}E_C,\qquad
 v_{u,\sigma}=\sigma\sqrt2 E_Ce_u
 \label{ii:ico:proyector-seis}
\end{equation}
realize \(\mathcal D_C\) as twelve distinct unit vectors in a
three-dimensional Euclidean space. There is an explicit isometry
\(T_C:\mathcal H_C\to\mathbb R^3\) for which
\(\iota_C(u,\sigma)=T_Cv_{u,\sigma}\) is a bijection onto
\(\Omega_{12}^{\rm ico}\), preserves incidence, and satisfies
\(\iota_C\jmath_C=-\iota_C\).
\end{proposition}
\end{samepage}
\begin{proof}
From \(C^2=5I_6\), one obtains \(E_C^2=E_C=E_C^*\); its trace is three.
The metric identities are
\[
 \langle v_{u,\sigma},v_{v,\tau}\rangle
 =\begin{cases}\sigma\tau,&u=v,\\
 \sigma\tau C_{uv}/\sqrt5,&u\ne v.
 \end{cases}
\]
Thus the twelve vectors are distinct, and their inner-product
\(1/\sqrt5\) incidence is exactly~\eqref{ii:ico:incidencia-firmada}.
To make the isometry explicit, let the columns of \(B\), in order
\(\mathcal U\), be
\[
 B=\frac1\nu
 \begin{pmatrix}
 0&-1&-\varphi_g&0&\varphi_g&1\\
 -1&-\varphi_g&0&1&0&-\varphi_g\\
 -\varphi_g&0&-1&-\varphi_g&-1&0
 \end{pmatrix}.
\]
Substitution of \(\varphi_g^2=\varphi_g+1\) gives
\(B^*B=I_6+C/\sqrt5=2E_C\). Hence
\(T_C=B/\sqrt2\big|_{\mathcal H_C}\) is an isometry, and
\(T_Cv_{u,\sigma}=\sigma Be_u\). The six columns of \(B\) and their
opposites are precisely the twelve vertices of
\eqref{ii:pent:eq:gravity-icosahedron-vertices}. The last identity
also proves compatibility of the involutions.
\end{proof}

The preceding dodecaphasic realization constructs the twelve vectors
directly from the projector \(P_3^{\mathrm{ico}}\) and the calendar positions.
The signed covering provides another coordinate expression of that
carrier. For the reader \(\pi_{\rm ico}\) in the descent theorem,
the relation between the two expressions is written
\begin{equation}
 X_{\rm TPK}^{\rm exc}\xrightarrow{\ r_\pm\ }
 \mathcal D_C\xrightarrow[\simeq]{\ \iota_C\ }
 \Omega_{12}^{\rm ico},\qquad
 \pi_{\rm ico}(x)=\sigma(x)Be_{u(x)},\quad
 r_\pm(x)=(u(x),\sigma(x)).
 \label{ii:ico:factorizacion-lector}
\end{equation}
Since \(\iota_C\) is a bijection, this factorization determines
\(r_\pm=\iota_C^{-1}\pi_{\rm ico}\). Thus \(r_\pm\) expresses
the reader being used in the signed chart: its introduction
does not add an independent hypothesis for constructing the dodecaphasic
carrier. The preceding correspondence table gives this expression
for its twelve positions. Properties (i)--(iv) retain their
own scope when the realization is identified with the quotient
of the full exceptional route. In particular, joint transport
of frame, projector, and incidence proves covariance between charts;
preservation of a fixed incidence by an active transformation
is a different identity. Orientation and memory remain
in the source state during both readings.
A signed permutation transporting \(C\) jointly transports
\(E_C\), its vectors, and incidence; this covariance allows the
charts to be compared while preserving each fiber's orientation.
